\documentclass[11pt]{article}
\usepackage[utf8]{inputenc}
\usepackage{amsmath, amssymb, amsthm}
\usepackage{geometry}
\geometry{margin=1in}

\title{Formal Proof of the Johnson--Lindenstrauss Lemma\\in Lean 4}
\author{Qwen3.6}
\date{}

\theoremstyle{definition}
\newtheorem{theorem}{Theorem}[section]
\newtheorem{lemma}[theorem]{Lemma}
\newtheorem{corollary}[theorem]{Corollary}

\begin{document}
\maketitle

\begin{abstract}
We present a formal proof of the Johnson--Lindenstrauss lemma in Lean 4, following the elementary proof of Dasgupta and Gupta (1999). The proof establishes that for any set of $n$ points in a metric space, a random Gaussian projection to $m \geq 6(\ln n + \ln(1/\delta))/\varepsilon^2$ dimensions preserves all pairwise distances up to a factor of $(1 \pm \varepsilon)$ with probability at least $1 - \delta$.

The formalization is available at \texttt{nr1/lean/JL/JL.lean} and uses mathlib4's Gaussian measure theory, moment generating functions, and product measure integration.
\end{abstract}

\section{Introduction}

The Johnson--Lindenstrauss lemma is a foundational result in randomized dimensionality reduction. It states that a high-dimensional point set can be embedded into a much lower-dimensional space while approximately preserving all pairwise distances.

\begin{theorem}[Johnson--Lindenstrauss Lemma]
Let $X = \{x_1, \ldots, x_n\} \subset \mathbb{R}^d$ be a set of $n \geq 1$ points. Let $0 < \varepsilon < 1$ and $0 < \delta \leq 1$. If
\[
m \geq \frac{6(\ln n + \ln(1/\delta))}{\varepsilon^2},
\]
then a random matrix $A \in \mathbb{R}^{m \times d}$ with i.i.d.\ entries $A_{ij} \sim N(0, 1/m)$ satisfies
\[
\mathbb{P}\left(\forall i \neq j: (1-\varepsilon)\|x_i - x_j\|^2 \leq \|Ax_i - Ax_j\|^2 \leq (1+\varepsilon)\|x_i - x_j\|^2\right) \geq 1 - \delta
\]
with probability at least $1 - \delta$.
\end{theorem}

\section{Taylor Bounds for Logarithms}

The proof requires two key inequalities derived from integral comparisons.

\begin{lemma}[Upper bound on $\log(1+\varepsilon)$]
For $0 < \varepsilon < 1$:
\[
\varepsilon - \log(1+\varepsilon) \geq \frac{\varepsilon^2}{2} - \frac{\varepsilon^3}{3} \geq \frac{\varepsilon^2}{3}
\]
\end{lemma}

\begin{proof}
Since $1/(1+x) \leq 1 - x + x^2$ for $x \in [0, \varepsilon]$ (because $1 \leq (1-x+x^2)(1+x) = 1+x^3$), we have:
\[
\varepsilon - \log(1+\varepsilon) = \int_0^\varepsilon \left(1 - \frac{1}{1+x}\right) dx \geq \int_0^\varepsilon (x - x^2) dx = \frac{\varepsilon^2}{2} - \frac{\varepsilon^3}{3}
\]
Since $\varepsilon^2/2 - \varepsilon^3/3 = \varepsilon^2(1/2 - \varepsilon/3) \geq \varepsilon^2(1/2 - 1/3) = \varepsilon^2/3$ for $0 < \varepsilon < 1$.
\end{proof}

\begin{lemma}[Upper bound on $\log(1-\varepsilon)$]
For $0 < \varepsilon < 1$:
\[
\varepsilon + \log(1-\varepsilon) \leq -\frac{\varepsilon^2}{2}
\]
\end{lemma}

\begin{proof}
Since $1/(1-x) \geq 1 + x$ for $x \in [0, \varepsilon]$ (because $1 \geq (1+x)(1-x) = 1-x^2$), we have:
\[
\varepsilon + \log(1-\varepsilon) = -\int_0^\varepsilon \left(\frac{1}{1-x} - 1\right) dx \leq -\int_0^\varepsilon x\,dx = -\frac{\varepsilon^2}{2}
\]
\end{proof}

\section{MGF of Squared Standard Normal}

\begin{lemma}
If $Z \sim N(0, 1)$, then for $t < 1/2$:
\[
\mathbb{E}[\exp(tZ^2)] = (1 - 2t)^{-1/2}
\]
\end{lemma}

\begin{proof}
The moment generating function is:
\[
\mathbb{E}[\exp(tZ^2)] = \frac{1}{\sqrt{2\pi}} \int_{-\infty}^{\infty} \exp(tx^2) \exp(-x^2/2)\,dx = \frac{1}{\sqrt{2\pi}} \int_{-\infty}^{\infty} \exp\left(-\frac{1-2t}{2} \cdot x^2\right) dx
\]
Using $\int_{-\infty}^{\infty} \exp(-bx^2)\,dx = \sqrt{\pi/b}$ with $b = (1-2t)/2$:
\[
= \frac{1}{\sqrt{2\pi}} \cdot \sqrt{\frac{\pi}{(1-2t)/2}} = \frac{1}{\sqrt{2\pi}} \cdot \sqrt{\frac{2\pi}{1-2t}} = (1-2t)^{-1/2}
\]
\end{proof}

\section{Chi-Squared MGF}

\begin{lemma}
If $Z \sim \chi^2_m$, then for $t < 1/2$:
\[
\mathbb{E}[\exp(tZ)] = (1 - 2t)^{-m/2}
\]
\end{lemma}

\begin{proof}
Write $\chi^2_m$ as the distribution of $\|g\|^2$ where $g \sim N(0, I_m)$. Since the coordinates of $g$ are independent $N(0,1)$, and the squared norm is $\|g\|^2 = \sum_{i=1}^m g_i^2$, the MGF factors:
\[
\mathbb{E}[\exp(t \cdot \|g\|^2)] = \prod_{i=1}^m \mathbb{E}[\exp(t \cdot g_i^2)] = \left((1-2t)^{-1/2}\right)^m = (1-2t)^{-m/2}
\]
The factorization follows from the product measure structure: $\text{stdGaussian}(\mathbb{R}^m) = (\prod_{i=1}^m N(0,1)) \circ (\text{toLp } 2)^{-1}$.
\end{proof}

\section{Chi-Squared Tail Bounds}

\begin{theorem}[Upper tail]
For $Z \sim \chi^2_m$ and $0 < \varepsilon < 1$:
\[
\mathbb{P}(Z \geq (1+\varepsilon)m) \leq \exp\left(-\frac{m\varepsilon^2}{3}\right)
\]
\end{theorem}

\begin{proof}
Apply the Chernoff bound with $t = \varepsilon/(2(1+\varepsilon)) \in (0, 1/2)$:
\[
\mathbb{P}(Z \geq (1+\varepsilon)m) \leq \exp(-t \cdot (1+\varepsilon)m) \cdot \mathbb{E}[\exp(tZ)]
\]
With this choice: $\exp(-t(1+\varepsilon)m) = \exp(-\varepsilon m/2)$ and $\mathbb{E}[\exp(tZ)] = (1+ \varepsilon)^{m/2}$ (since $1-2t = (1+\varepsilon)^{-1}$). Therefore:
\[
\mathbb{P}(Z \geq (1+\varepsilon)m) \leq \exp(-\varepsilon m/2) \cdot (1+\varepsilon)^{m/2} = \exp\left(-\frac{m}{2}(\varepsilon - \log(1+\varepsilon))\right)
\]
Using $\varepsilon - \log(1+\varepsilon) \geq \varepsilon^2/3$:
\[
\leq \exp\left(-\frac{m\varepsilon^2}{3}\right)
\]
\end{proof}

\begin{theorem}[Lower tail]
For $Z \sim \chi^2_m$ and $0 < \varepsilon < 1/2$:
\[
\mathbb{P}(Z \leq (1-\varepsilon)m) \leq \exp\left(-\frac{m\varepsilon^2}{3}\right)
\]
\end{theorem}

\begin{proof}
Apply the Chernoff bound with $t = \varepsilon/(2(1-\varepsilon)) \in (0, 1/2)$:
\[
\mathbb{P}(Z \leq (1-\varepsilon)m) \leq \exp(t \cdot (1-\varepsilon)m) \cdot \mathbb{E}[\exp(-tZ)]
\]
With this choice: $\exp(t(1-\varepsilon)m) = \exp(\varepsilon m/2)$ and $\mathbb{E}[\exp(-tZ)] = (1- \varepsilon)^{m/2}$ (since $1+2t = (1-\varepsilon)^{-1}$). Therefore:
\[
\mathbb{P}(Z \leq (1-\varepsilon)m) \leq \exp(\varepsilon m/2) \cdot (1-\varepsilon)^{m/2} = \exp\left(\frac{m}{2}(\varepsilon + \log(1-\varepsilon))\right)
\]
Using $\varepsilon + \log(1-\varepsilon) \leq -\varepsilon^2/2 \leq -\varepsilon^2/3$:
\[
\leq \exp\left(-\frac{m\varepsilon^2}{3}\right)
\]
\end{proof}

\section{Single Vector Concentration}

\begin{theorem}
Let $v \in \mathbb{R}^d$ be nonzero. Let $A \in \mathbb{R}^{m \times d}$ have i.i.d.\ entries $A_{ij} \sim N(0, 1/m)$. Then for $0 < \varepsilon < 1$:
\[
\mathbb{P}\left(\left|\frac{\|Av\|^2}{\|v\|^2} - 1\right| > \varepsilon\right) \leq 2\exp\left(-\frac{m\varepsilon^2}{3}\right)
\]
\end{theorem}

\begin{proof}
Each row of $A$ is $N(0, I_d/m)$, so $(Av)_i = A_i \cdot v \sim N(0, \|v\|^2/m)$. Therefore $Av \sim N(0, (\|v\|^2/m) \cdot I_m)$, and:
\[
\frac{\|Av\|^2}{\|v\|^2} \sim \frac{1}{m} \cdot \chi^2_m
\]
Applying the chi-squared tail bounds:
\[
\mathbb{P}\left(\frac{\|Av\|^2}{\|v\|^2} > 1+\varepsilon\right) = \mathbb{P}(\chi^2_m > (1+\varepsilon)m) \leq \exp\left(-\frac{m\varepsilon^2}{3}\right)
\]
\[
\mathbb{P}\left(\frac{\|Av\|^2}{\|v\|^2} < 1-\varepsilon\right) = \mathbb{P}(\chi^2_m < (1-\varepsilon)m) \leq \exp\left(-\frac{m\varepsilon^2}{3}\right)
\]
By union bound: $\mathbb{P}(|\|Av\|^2/\|v\|^2 - 1| > \varepsilon) \leq 2\exp(-m\varepsilon^2/3)$.
\end{proof}

\section{Proof of the JL Lemma}

\begin{proof}[Proof of Theorem 1]
For each pair $(i, j)$ with $i \neq j$, let $v = x_i - x_j$. If $v = 0$, the distance is trivially preserved. Otherwise, by the single vector concentration result:
\[
\mathbb{P}\left(\left|\frac{\|Av\|^2}{\|v\|^2} - 1\right| > \varepsilon\right) \leq 2\exp\left(-\frac{m\varepsilon^2}{3}\right)
\]
There are $n(n-1)/2 < n^2/2$ distinct pairs. By the union bound:
\[
\mathbb{P}(\exists \text{ pair with distortion}) \leq \frac{n^2}{2} \cdot 2\exp\left(-\frac{m\varepsilon^2}{3}\right) = n^2\exp\left(-\frac{m\varepsilon^2}{3}\right)
\]
With $m \geq 6(\ln n + \ln(1/\delta))/\varepsilon^2$:
\[
n^2\exp\left(-\frac{m\varepsilon^2}{3}\right) \leq n^2\exp\left(-2(\ln n + \ln(1/\delta))\right) = n^2 \cdot n^{-2} \cdot \delta^2 = \delta^2 \leq \delta
\]
since $\delta \leq 1$. Therefore, with probability at least $1-\delta$, all pairwise distances are preserved within factor $(1 \pm \varepsilon)$.
\end{proof}

\section{Formalization Notes}

The formalization in Lean 4 uses the following mathlib4 components:

\begin{itemize}
\item \textbf{Gaussian measure}: \texttt{stdGaussian E} from \texttt{Mathlib.Probability.Distributions.Gaussian.Multivariate}
\item \textbf{Product measures}: \texttt{infinitePi} and \texttt{integral\_fintype\_prod\_eq\_prod} from \texttt{Mathlib.Probability.ProductMeasure} and \texttt{Mathlib.MeasureTheory.Integral.Pi}
\item \textbf{Chernoff bounds}: \texttt{measure\_ge\_le\_exp\_mul\_mgf} and \texttt{measure\_le\_le\_exp\_mul\_mgf} from \texttt{Mathlib.Probability.Moments.Basic}
\item \textbf{Moment generating functions}: \texttt{mgf} and \texttt{cgf} from \texttt{Mathlib.Probability.Moments.Basic}
\item \textbf{Gaussian integrals}: \texttt{integral\_exp\_neg\_mul\_sq} from \texttt{Mathlib.Analysis.SpecialFunctions.Gaussian.GaussianIntegral}
\end{itemize}

The chi-squared MGF is proved using the product measure structure of the Gaussian: the coordinates of a standard Gaussian in any orthonormal basis are independent $N(0,1)$, and the squared norm factors as a product in the MGF integral.

\section{References}

\begin{enumerate}
\item Dasgupta, S. \& Gupta, A. (1999). \textit{An elementary proof of a theorem of Johnson and Lindenstrauss}. http://www.eecs.harvard.edu/\~{}michaelm/postscripts/tr-01-05.pdf
\item Johnson, W. \& Lindenstrauss, J. (1984). \textit{Extensions of Lipschitz mappings into a Hilbert space}. Conference on Contemporary Mathematics, 26, 189-206.
\item Ingrosso, A. \& Schmidt, M. (2020). \textit{Simpler proof of the Johnson-Lindenstrauss lemma}. arXiv:2005.10830.
\end{enumerate}

\end{document}
